% !Mode:: "TeX:UTF-8"
% !TEX program  = xelatex

\section{实验设计}
\label{sec:experiment}

\subsection{数据集}

实验包含两类数据设置，分别用于评估同分布性能和分布外泛化。
下文将其记为 \textbf{合成混合数据集} 和 \textbf{LIB 基准数据集}。

\paragraph{合成混合数据集。}
主要实验使用前述工作整理并公开的数据划分。
TSP 与 CVRP 各自包含 10000 个训练实例、1000 个验证实例和 1000 个测试实例\cite{gao2025nss}。
两类问题的规模范围均为 50 到 499，实例分布由均匀分布和高斯混合分布组成。

\paragraph{LIB 基准数据集。}
为了检验方法在真实基准集上的泛化能力，
本文进一步在 TSPLIB 与 CVRPLIB 组成的 LIB 基准数据集上评测验证阶段选出的最佳选择模型。
其中 TSPLIB\cite{reinelt1991tsplib} 共 49 个实例，CVRPLIB\cite{uchoa2017cvrplib} 共 100 个实例。

\subsection{评价指标}
\begin{itemize}
  \item \textbf{mean length}：选中方法在测试集上的平均路径长度，数值越小越好；
  \item \textbf{top-1 accuracy}：选择器选中了该实例上实际最优方法（即与 Oracle 一致）的比例，数值越大越好。
\end{itemize}

\subsection{基线方法}
\begin{itemize}
  \item \textbf{Oracle}：对每个实例选择其实际最优方法（即事后最优），作为选择性能的理论上界；
  \item \textbf{SingleBest}：对每个问题固定选择测试集上平均表现最好的单一方法。注意该基线使用了测试集的全局统计信息来确定固定方法，因此不是一个可在线部署的策略，而是作为固定选择的性能上界；
  \item \textbf{Neural-LinUCB}：本文的主方法，采用深度表征与浅层线性探索解耦\cite{xu2022neurallinucb}。
\end{itemize}

\subsection{训练设置}

Neural-LinUCB 的完整超参数见表 \ref{tab:default-hparams}。模型每次只根据被选动作接收反馈并更新统计量。
